%! AP Statistics — Unit 4, Lesson 4.11 — Group Activity KEY
\documentclass[10pt]{article}
\usepackage{apstats-article}
\usepackage{apstats-key}

\begin{document}

\pageheader{Unit 4, Lesson 4.11}{Group Activity KEY: Binomial Parameters in Context}
\namepartnerperiod

\vspace{0.1in}

\begin{scenariobox}[Directions]{sky}
\small
Work through all three tiers. Show all formula setup before computing. Write
complete-sentence interpretations where asked --- the AP exam requires this.
\end{scenariobox}

\vspace{0.1in}

% ── TIER R ────────────────────────────────────────────────────────────────────
\begin{tcolorbox}[colback=goldbg, colframe=goldacc, boxrule=1pt, arc=2mm,
    title={\bfseries Tier R \quad Remediate: Computing Mean and SD},
    fonttitle=\small]
\small

For each binomial distribution, compute $\mu_X$ and $\sigma_X$, then write a
one-sentence interpretation of $\mu_X$.

\vspace{10pt}
\textbf{(a)} $X \sim B(25, 0.40)$

\vspace{4pt}
$\mu_X = np = $ \ans{25(0.40)} $= $ \ans{10}

$\sigma_X = \sqrt{np(1-p)} = \sqrt{{\color{keyred}\mathbf{25(0.40)(0.60)}}} \approx$ \ans{2.45}

Interpretation of $\mu_X$: \ansline{Over many repetitions,}\\[1pt]\ansline{the average number of successes would be 10.}

\vspace{12pt}
\textbf{(b)} $X \sim B(100, 0.08)$

\vspace{4pt}
$\mu_X = $ \ans{100(0.08)} $= $ \ans{8}

$\sigma_X = \sqrt{{\color{keyred}\mathbf{100(0.08)(0.92)}}} \approx$ \ans{2.71}

Interpretation of $\mu_X$: \ansline{Over many repetitions,}\\[1pt]\ansline{the average number of successes would be 8.}
\vspace{4pt}
\end{tcolorbox}

\vspace{0.1in}

% ── TIER A ────────────────────────────────────────────────────────────────────
\begin{tcolorbox}[colback=sky, colframe=navy, boxrule=1pt, arc=2mm,
    title={\bfseries Tier A \quad Approaching Proficiency: A Full Binomial Analysis},
    fonttitle=\small]
\small

\textbf{Context:} A call center reports that 22\% of customer calls are resolved
on the first attempt. On a busy afternoon, 80 calls come in. Assume $X \sim B(80, 0.22)$.

\vspace{8pt}
(a) Calculate $\mu_X$ and $\sigma_X$. Show formula setup.

\vspace{4pt}
$\mu_X = $ \ans{80(0.22)} $= $ \ans{17.6}

$\sigma_X = $ \ans{$\sqrt{80(0.22)(0.78)} = \sqrt{13.728}$} $\approx$ \ans{3.71}

\vspace{8pt}
(b) Interpret $\mu_X$ in context.

\vspace{4pt}
\ansline{Over many afternoons of 80 calls, the center would average}\\[1pt]\ansline{about 17.6 calls resolved on the first attempt per afternoon.}

\vspace{8pt}
(c) Find $P(X \le 15)$. Write the calculator notation and the answer.

\vspace{4pt}
$P(X \le 15) = \texttt{binomcdf(}{\color{keyred}\mathbf{80}}\texttt{,}\;{\color{keyred}\mathbf{0.22}}\texttt{,}\;15\texttt{)} \approx $ \ans{0.2752}

\vspace{8pt}
(d) Find $P(X \ge 20)$. Write the inequality rewrite, then the calculator notation.

\vspace{4pt}
$P(X \ge 20) = 1 - P(X \le {\color{keyred}\mathbf{19}}) = 1 - \texttt{binomcdf(}{\color{keyred}\mathbf{80}}\texttt{,}\;{\color{keyred}\mathbf{0.22}}\texttt{,}\;{\color{keyred}\mathbf{19}}\texttt{)} \approx $ \ans{0.2695}

\vspace{8pt}
(e) Is 20 or more resolved calls on the first attempt \emph{unusually high}?
Use $\mu_X$ and $\sigma_X$ to support your answer.

\vspace{4pt}
\ansline{No. $\mu_X + 2\sigma_X = 17.6 + 2(3.71) \approx 25.0$; 20 is only about}\\[1pt]\ansline{$0.65\sigma_X$ above the mean, so it is not unusually high.}
\vspace{4pt}
\end{tcolorbox}

\vspace{0.1in}

% ── TIER E ────────────────────────────────────────────────────────────────────
\begin{tcolorbox}[colback=greenbg, colframe=greenacc, boxrule=1pt, arc=2mm,
    title={\bfseries Tier E \quad Extension: Comparing Two Binomials},
    fonttitle=\small]
\small

\textbf{Context:} Two quality-control processes are being evaluated.
\begin{itemize}[itemsep=3pt]
  \item Process A: $n = 200$ items inspected, defect rate $p = 0.05$.
  \item Process B: $n = 100$ items inspected, defect rate $p = 0.10$.
\end{itemize}

\vspace{6pt}
(a) Compute $\mu_X$ and $\sigma_X$ for both processes.

\vspace{4pt}
\renewcommand{\arraystretch}{1.8}
\begin{tabularx}{\linewidth}{@{} >{\bfseries}p{2cm} X X @{}}
 & Process A & Process B \\
\hline
$\mu_X$ & \ans{$200(0.05) = 10$} & \ans{$100(0.10) = 10$} \\
$\sigma_X$ & \ans{$\sqrt{200(0.05)(0.95)} \approx 3.08$} & \ans{$\sqrt{100(0.10)(0.90)} = 3.00$} \\
\end{tabularx}

\vspace{8pt}
(b) Both processes have the same mean. Explain why in terms of the formulas.

\vspace{4pt}
\ansline{$\mu_X = np$ and $200(0.05) = 100(0.10) = 10$:}\\[1pt]\ansline{more trials at a lower rate give the same expected count.}

\vspace{8pt}
(c) Which process has more variability? What does this mean practically for the
quality-control engineer?

\vspace{4pt}
\ansline{Process A ($\sigma_X \approx 3.08$ vs.\ $3.00$), slightly more.}\\[1pt]\ansline{Its defect counts fluctuate a bit more from batch to batch.}

\vspace{8pt}
(d) Challenge: For a fixed $n$, what value of $p$ maximizes $\sigma_X$?
Show a brief algebraic or numerical argument.

\vspace{4pt}
\ansline{$\sigma_X = \sqrt{np(1-p)}$ is largest when $p(1-p)$ is largest.}\\[1pt]\ansline{$p(1-p)$: 0.09 at $p=0.1$, 0.21 at 0.3, 0.25 at 0.5, 0.21 at 0.7.}\\[1pt]\ansline{Symmetric peak at $p = 0.5$, so $\sigma_X$ is maximized at $p = 0.5$.}
\vspace{4pt}
\end{tcolorbox}

\end{document}
